% !TeX root = ../main.tex
\begin{frame}{03 创新思路——验证思路与对照}
  \begin{SlideCanvas}
    \BodyAt{55}{84}{850}{在统一材料、尺度与制造约束下，比较长度误差、径向力曲线和局部应变，逐层检验拟议设计的可行范围。}
    \PointAt{55}{153}{175}{验证层级}
    \PointAt{267}{153}{268}{加载与边界}
    \PointAt{595}{153}{305}{关键指标}
    \MainRuleAt{191}
    \PointAt{55}{214}{175}{单胞}
    \DetailAt{267}{214}{268}{周向拉伸，轴向自由，周期性边界条件}
    \DetailAt{595}{214}{305}{泊松比、应力响应及目标曲线拟合误差}
    \RuleAt{270}
    \PointAt{55}{293}{175}{短段圆管}
    \DetailAt{267}{293}{268}{径向扩张与规定压缩，保留轴向自由度}
    \DetailAt{595}{293}{305}{$L(D)$、$F_r(D)$\\端部位移、扭转与局部应变}
    \RuleAt{349}
    \PointAt{55}{372}{175}{血管接触}
    \DetailAt{267}{372}{268}{匹配管腔恢复程度，采用对应部署本构}
    \DetailAt{595}{372}{305}{贴壁间隙与血管壁应变，必要时增加血流评价}
    \CaptionAt{55}{446}{850}{对照：常规构型、文献零泊松比构型、变形控制、联合设计。统一材料量与最低承载要求。}
    \Sources{拟议验证与对照方案。方法：\RefGroup；器件参考：\RefWang、\RefYasmin。}
  \end{SlideCanvas}
  \BICConclusionBand{在相同长度误差水平下，检验\BICBandEmph{力曲线定制的收益与应变代价}}
  \note{本页合并原有对照与验证两页，移除粗略示意，用加载、边界与指标的对应表具体说明研究设计。四类对照中，变形控制和联合设计共享设计域及材料、制造、应变和最低承载条件。长度误差取直径区间最大绝对相对偏差；力曲线误差按同一加载分支以参考力归一；同时报告局部应变和端部变形。应扫描目标与约束形成可行范围，避免只挑选成功解。血管接触要匹配管腔恢复程度。实际永久展开金属需要弹塑性，自膨胀NiTi需要相变超弹性及压握释放路径。}
\end{frame}
